% META: {"id": "TCOMPL-ATT-CO-026", "chapitre": "TCOMPL-TEMPS-ATTENTE", "type_objet": "corrige", "exercice_ref": "TCOMPL-ATT-EX-026", "capacites_codes": ["C2"], "sources_inspiration": [], "status": "approved"}
% BEGIN-VERIFY
% from sympy import *
% a, b, c = symbols('a b c', positive=True)
% # AM-GM pour 3 variables: (a+b+c)/3 >= (abc)^(1/3)
% # par concavite de ln
% lhs = ln((a+b+c)/3)
% rhs = (ln(a) + ln(b) + ln(c))/3
% # numerique
% assert lhs.subs({a:1, b:2, c:4}) == ln(Rational(7,3))
% assert simplify(rhs.subs({a:1, b:2, c:4}) - ln(2)) == 0
% assert ln(Rational(7,3)) > ln(2)
% END-VERIFY
\begin{corrige}{TCOMPL-ATT-EX-026}

La fonction $\ln$ est concave. Par Jensen avec $t_1=t_2=t_3=\dfrac{1}{3}$ :
\[ \ln\!\left(\frac{a+b+c}{3}\right) \geq \frac{\ln a + \ln b + \ln c}{3} = \ln\sqrt[3]{abc}. \]

\commentaireMarge{Jensen avec 3 points et poids $1/3$ chacun.}

Comme $\ln$ est croissante : $\dfrac{a+b+c}{3} \geq \sqrt[3]{abc}$.

L'egalite a lieu si et seulement si $a = b = c$.

\end{corrige}
